\section{ТЕОРЕТИЧЕСКИЕ ОСНОВЫ И АНАЛИТИЧЕСКИЙ ОБЗОР ПОДХОДОВ К УПРАВЛЕНИЮ ИНФРАСТРУКТУРОЙ}

Инфраструктура современной программной системы включает вычислительные узлы, сетевые настройки, контейнеры, системные сервисы, средства наблюдаемости и прикладные workload. При росте числа команд, окружений и управляемых хостов поддержание согласованного состояния перестает быть задачей разовой настройки и превращается в непрерывный эксплуатационный процесс.

В рамках данной работы инфраструктура рассматривается не только как набор ресурсов, которые необходимо создать, но и как изменяемая система, состояние которой должно постоянно сопоставляться с архитектурным описанием. Такой взгляд близок к модели управления Kubernetes: пользователь описывает желаемое состояние объектов, а контроллеры наблюдают фактическое состояние и выполняют действия для приближения системы к целевому виду \cite{kubernetes_controllers,kubernetes_objects}. Отличие рассматриваемой платформы состоит в том, что она ориентирована не на внутренние объекты Kubernetes-кластера, а на более общий контур управления архитектурных манифестов.

В данной главе рассматриваются ключевые понятия, необходимые для постановки задачи: desired state, actual state, drift, reconcile, Infrastructure-as-Code и Architecture-as-Code. Далее проводится анализ существующих решений и формулируются ограничения, которые мотивируют разработку собственной платформы.

\subsection{Проблема управления инфраструктурой в распределённых системах}

Распределенная система редко существует в единственном экземпляре. Как правило она разворачивается в нескольких средах: разработки, тестовой, preproduction и production. Каждая среда содержит собственный набор хостов, контейнеров, вспомогательных сервисов, источников телеметрии и конфигураций. Даже если первоначально такие среды настроены одинаково, со временем между ними возникают расхождения.

Причины расхождений разнообразны. Часть изменений вносится вручную в ходе оперативного устранения инцидентов. Часть появляется из-за неполных или устаревших инструкций развертывания. Часть связана с особенностями конкретного окружения: отличающимися версиями пакетов, сетевыми ограничениями, частично выполненными rollout-процедурами. В результате фактическое состояние инфраструктуры постепенно отходит от того состояния, которое команда считает правильным.

Для SRE- и DevOps-команд это приводит к росту операционных затрат. Чем больше инфраструктурных объектов находится в зоне ответственности команды, тем дороже становится ручная верификация. Инженеры вынуждены поддерживать инвентарные списки, сверять конфигурации, запускать отдельные playbook, анализировать журналы выполнения и самостоятельно решать, какие изменения допустимо применить. Такой подход может быть приемлемым для небольшого числа хостов, но плохо масштабируется при увеличении числа окружений и команд.

Отдельный риск связан с тем, что инфраструктурное состояние часто описано в разных местах. Архитектурная документация может существовать отдельно от inventory, inventory отдельно от конфигурационных playbook, а фактическая информация о запущенных workload извлекаться из средств мониторинга или вручную с хостов. Например, сведения о контейнерах могут поступать из cAdvisor, а сведения о хостах и их атрибутах --- из систем инвентаризации уровня NetBox \cite{cadvisor_docs,netbox_docs}. Пока эти источники не сведены в единый управляемый контур, команда не получает надежного ответа на вопрос: соответствует ли инфраструктура тому, что было спроектировано.

Следовательно, задача управления инфраструктурой в распределенных системах должна рассматриваться как задача непрерывного сопоставления архитектурного описания и фактического состояния. Для этого необходимы: формализованный источник целевого состояния, механизм получения актуального состояния, компонент сравнения, правила безопасной обработки неполных данных и механизм приведения системы к целевому виду.

\subsection{Desired state, actual state, drift и reconcile как базовая модель управления}

В основе автоматизированного управления инфраструктурой лежит различение желаемого и фактического состояния. Desired state описывает, каким должно быть состояние системы: какие хосты находятся под управлением, какие workload должны быть запущены, какие параметры необходимы для их развертывания, какие атрибуты характеризуют хосты и их назначение. Actual state, напротив, фиксирует состояние, наблюдаемое в текущий момент: какие хосты доступны, какие workload реально присутствуют, какие источники данных успешно отработали и где возникли ошибки.

Такое разделение позволяет перейти от императивной модели ``выполнить набор команд'' к декларативной модели ``поддерживать заданное состояние''. В императивной модели результат зависит от того, когда и как была выполнена команда. В декларативной модели главным артефактом становится целевое описание, а управляющая система должна регулярно проверять, достигнуто ли оно. Именно эта идея лежит в основе Kubernetes-контроллеров: контроллер наблюдает состояние ресурсов и выполняет действия, необходимые для приближения текущего состояния к желаемому \cite{kubernetes_controllers}.

Drift представляет собой обнаруженное расхождение между desired state и actual state. В простейшем случае drift означает, что компонент, обязательный в desired state, отсутствует в actual state. Более сложные виды drift могут включать несоответствие версии, конфигурации, режима запуска, сетевых параметров.

Reconcile обозначает процесс устранения drift, то есть выполнение действий, направленных на приведение фактического состояния к желаемому. Это управляемое действие, основанное на конкретном результате сравнения и ограниченное заранее определенным набором операторов. Такой подход снижает риск несанкционированных изменений: платформа не выполняет произвольный код, а выбирает один из поддержанных сценариев восстановления.

Для корректной реализации данной модели важна семантика неполных данных. Если источник actual state недоступен или вернул частичный результат, отсутствие информации не всегда можно трактовать как drift. Например, если система не смогла получить данные от источника наблюдения, она не должна автоматически устанавливать все отсутствующие в ответе компоненты. Без такого ограничения платформа может начать выполнять лишние reconcile-действия из-за сетевой ошибки или временной недоступности источника.

Поэтому безопасная control-loop модель должна различать три ситуации:
\begin{enumerate}
    \item компонент присутствует в desired state и подтвержденно отсутствует в actual state;
    \item компонент присутствует в desired state, но actual state неполон;
    \item desired state не требует компонента или компонент выключен.
\end{enumerate}

Только первая ситуация является основанием для reconcile. Вторая ситуация должна приводить к предупреждению, partial-статусу или пропуску части проверки. Третья не требует действий. Такая семантика особенно важна для инфраструктурной платформы, так как ошибочные автоматические изменения потенциально затрагивают рабочие хосты.

В результате desired state, actual state, drift и reconcile образуют базовую модель управления. Она позволяет описать инфраструктурный процесс как повторяемый цикл: получить целевое состояние, получить фактическое состояние, сравнить их, зафиксировать расхождения, выполнить допустимые действия восстановления и повторить проверку.

\subsection{Подходы Infrastructure-as-Code и Architecture-as-Code}

Infrastructure-as-Code является устоявшимся подходом, при котором инфраструктурные ресурсы описываются в машинно-читаемом виде. Такие описания могут храниться в системе контроля версий, проходить ревью и применяться автоматизированными инструментами. Преимущество Infrastructure-as-Code состоит в воспроизводимости: конфигурация инфраструктуры становится артефактом разработки, а не набором ручных действий конкретного инженера.

Классические инструменты Infrastructure-as-Code, такие как Terraform и Pulumi, ориентированы на описание ресурсов облаков, сетей, баз данных, Kubernetes-объектов и других инфраструктурных сущностей \cite{terraform_docs,pulumi_iac}. Их сильная сторона заключается в построении плана изменений, управлении зависимостями и повторяемом применении конфигурации.

Architecture-as-Code решает другую задачу. Если Infrastructure-as-Code описывает преимущественно инфраструктурные ресурсы и способы их создания, то Architecture-as-Code фиксирует архитектурное представление системы: компоненты, контейнеры, связи, deployment-узлы, среду размещения и дополнительные свойства.

Практическая ценность Architecture-as-Code проявляется в том, что архитектурная модель может стать источником desired state. Если deployment-описание содержит хосты, их атрибуты и размещенные на них workload, платформа может извлечь из него формализованное целевое состояние.

Для такого подхода подходит C4-модель, поскольку она задает несколько уровней описания программной системы: контекст, контейнеры, компоненты и код \cite{c4_model}. На уровне deployment-описания можно зафиксировать конкретные узлы и экземпляры контейнеров, то есть перейти от абстрактной архитектуры к инфраструктурно значимым данным. Structurizr DSL и экспортируемое JSON-представление дают техническую основу для такого машинного разбора \cite{structurizr_dsl}.

При этом Architecture-as-Code не заменяет Infrastructure-as-Code полностью. Эти подходы решают взаимодополняющие задачи. Infrastructure-as-Code отвечает за описание и создание ресурсов, а Architecture-as-Code позволяет связать архитектурное намерение с последующей проверкой фактической инфраструктуры.

\subsection{Анализ существующих решений}

Для обоснования необходимости собственной платформы был проведен анализ существующих решений. Сравнение выполнялось с точки зрения следующих критериев: наличие декларативного desired state, поддержка непрерывной реконсиляции, применимость вне Kubernetes, работа с host-level workload, сложность внедрения и соответствие идее Architecture-as-Code.

Ansible является одним из наиболее распространенных инструментов управления конфигурациями. Его важное преимущество заключается в agentless-подходе: для управления Linux-хостом обычно достаточно SSH-доступа. Это упрощает внедрение и снижает требования к целевым узлам.

Однако Ansible сам по себе не является непрерывным control plane. Playbook запускается как задача, выполняет набор операций и завершает работу. Если после этого на хосте произойдет drift, инструмент не обнаружит и не устранит его без внешней оркестрации. Кроме того, при росте числа хостов и групп inventory может стать сложным для сопровождения.

Terraform и Pulumi решают задачу декларативного описания инфраструктуры, но делают это преимущественно на уровне ресурсов и провайдеров. Они хорошо подходят для создания облачных объектов, сетей, managed-сервисов, Kubernetes-ресурсов и связанных зависимостей. Важное преимущество Terraform состоит в понятной модели plan/apply, а Pulumi --- в возможности использовать языки программирования для построения абстракций. Тем не менее оба подхода сохраняют характер запуска операции применения и не являются универсальным механизмом постоянного контроля host-level workload.

Kubernetes демонстрирует наиболее близкую к целевой платформе идею. В Kubernetes есть декларативные объекты, API-сервер, контроллеры и цикл согласования, который постоянно стремится привести observed state к desired state \cite{kubernetes_controllers,kubernetes_objects}. Операторы расширяют эту модель на доменные сущности \cite{kubernetes_operator_pattern}.

Ограничение Kubernetes-подхода состоит в его границах применимости. Если управляемая среда уже является Kubernetes-кластером, операторы и GitOps-инструменты дают сильный механизм реконсиляции. Но если необходимо управлять workload на VM, bare-metal-хостах или смешанной инфраструктуре, приходится либо переносить объекты в Kubernetes-модель, либо строить адаптеры.

Crossplane развивает идею Kubernetes как универсального control plane для внешних ресурсов. Через CRD и composition-механизм он позволяет описывать высокоуровневые инфраструктурные абстракции и управлять внешними объектами. Такой подход технологически мощен, но требует Kubernetes-кластера и компетенций в его расширении. Для задачи данной работы это важный ориентир, но не прямой заменитель: проектируемая платформа должна быть проще и фокусироваться на связке Architecture-as-Code, inventory, drift detection и reconcile.

Argo CD и Flux представляют GitOps-подход. Они автоматически синхронизируют Kubernetes-кластеры с конфигурациями, хранящимися в Git, и хорошо решают задачу drift для Kubernetes-ресурсов. Их ограничение связано с тем, что Git становится центральным runtime-источником изменений, а область управления остается в основном кластерной.

ONAP демонстрирует противоположный край спектра: это крупная платформа автоматизации жизненного цикла сетевых сервисов. Она содержит богатые подсистемы проектирования, inventory, политик, аналитики и closed-loop automation. В то же время ONAP ориентирована на телеком-домен и требует значительных ресурсов внедрения и эксплуатации. Для рассматриваемого проекта такой уровень комплексности является избыточным.

\begin{table}
    \centering
    \caption{Сравнение подходов и инструментов управления инфраструктурой}
    \label{tab:infra-tools-comparison}
    \footnotesize
    \begin{tabularx}{\textwidth}{|p{0.18\textwidth}|X|X|}\hline
    Инструмент & Сильные стороны & Ограничения относительно задачи работы \\ \hline
    Ansible & Agentless-модель, широкая экосистема модулей и ролей, удобство разового применения конфигураций по SSH \cite{ansible_docs}. & Job-based выполнение, отсутствие встроенной непрерывной реконсиляции, рост сложности inventory и playbook в крупных системах. \\ \hline
    Terraform & Декларативное описание ресурсов, план изменений, развитая экосистема провайдеров \cite{terraform_docs}. & Ориентация на provision инфраструктурных ресурсов и state-файл; host-level workload и непрерывная реконсиляция не являются основной моделью. \\ \hline
    Pulumi & Infrastructure-as-Code на языках программирования, возможность строить абстракции и тестировать код конфигурации \cite{pulumi_iac}. & Сохраняет job-based характер применения, требует навыков разработки и отдельного управления состоянием. \\ \hline
    Kubernetes Operators & Нативная control-loop модель, CRD и операторы для доменных объектов \cite{kubernetes_operator_pattern}. & Сильно привязаны к Kubernetes API и кластерной модели; bare-metal/VM-хосты и внешние workload требуют дополнительных адаптеров. \\ \hline
    Crossplane & Расширяет Kubernetes до control plane для внешних ресурсов, поддерживает композиции и непрерывное согласование \cite{crossplane_docs}. & Требует Kubernetes как базовую платформу, сложен для команд без Kubernetes-экспертизы, не фокусируется на host-level workload. \\ \hline
    Argo CD и Flux & GitOps-синхронизация Kubernetes-ресурсов, наблюдение drift в кластере, интеграция с Git \cite{argocd_docs,flux_docs}. & Ориентированы на Kubernetes-манифесты и Git как центральный источник; внешние хосты и архитектурные deployment-модели покрываются косвенно. \\ \hline
    ONAP & Комплексная автоматизация жизненного цикла сетевых сервисов, развитая модель inventory, политики и closed-loop automation \cite{onap_docs}. & Очень высокая сложность, телеком-ориентация, значительный операционный overhead для средних и небольших команд. \\ \hline
    \end{tabularx}
    \normalsize
\end{table}

Проведенный анализ показывает, что существующие инструменты закрывают отдельные части задачи, но не дают одновременно легковесной, domain-focused платформы, где Architecture-as-Code используется как источник desired state, inventory предоставляет actual state, отдельный сервис обнаруживает drift, а Reconciler выполняет набор действий восстановления.

\subsection{Ограничения существующих решений и мотивация собственной платформы}

Сравнение существующих решений позволяет выделить несколько повторяющихся ограничений. Первое ограничение связано с разрывом между архитектурным описанием и фактическим управлением инфраструктурой. Архитектура часто существует в виде диаграмм или текстовой документации, но эти артефакты не используются управляющими сервисами напрямую.

Второе ограничение состоит в job-based характере многих инструментов. Ansible, Terraform и Pulumi хорошо подходят для запуска конкретной операции, но для непрерывной реконсиляции требуется внешний цикл: расписание, pipeline, отдельный orchestrator или ручная инициатива инженера.

Третье ограничение относится к области применимости Kubernetes-ориентированных решений. Kubernetes Operators, Crossplane, Argo CD и Flux обладают сильной control-loop моделью, но их естественная среда --- Kubernetes API и Kubernetes-ресурсы. Для управления обычными хостами, Docker workload, агентами наблюдаемости или смешанной инфраструктурой их применение может оказаться слишком тяжелым.

Четвертое ограничение связано со сложностью крупных платформ. Системы уровня ONAP способны решать широкий набор задач автоматизации, но требуют значительного операционного overhead. Для команд, которым нужно решить конкретную проблему контроля drift и восстановления workload, внедрение такой платформы может быть несоразмерно сложным.

Мотивация собственной платформы заключается в поиске промежуточного решения. С одной стороны, оно должно использовать сильную идею Kubernetes-подобного control loop. С другой стороны, оно не должно требовать переноса всей инфраструктуры в Kubernetes или внедрения крупной enterprise-платформы. Такое решение должно быть понятным, расширяемым и пригодным для постепенного развития.

Для этого в проекте выбрана сервисная декомпозиция:
\begin{enumerate}
    \item ParserSvc отвечает за получение desired state из Architecture-as-Code-описания;
    \item InventorySvc отвечает за формирование actual state;
    \item DriftDetectorSvc сравнивает desired и actual state, обнаруживает drift и формирует reconcile-команды;
    \item ReconcilerSvc выполняет допустимые действия восстановления через операторную модель.
\end{enumerate}

Такое разделение позволяет независимо развивать источники desired state, источники actual state, набор detector-ов и набор reconcile-операторов. В MVP это выражено минимально: поддержаны Structurizr JSON как входной формат, cAdvisor\footnote{cAdvisor --- инструмент сбора и экспорта сведений о контейнерах и характеристиках их выполнения на хосте \cite{cadvisor_docs}.} как источник actual state и два workload-компонента: \texttt{node\_exporter} и \texttt{cadvisor}. Однако сама архитектура допускает расширение: добавление NetBox\footnote{NetBox --- система моделирования и документирования инфраструктуры, применяемая как source of truth для сетевых, серверных и связанных инфраструктурных ресурсов \cite{netbox_docs}.} как источника inventory, новых detector-ов, persistence-слоя, API истории drift и более сложных reconcile-сценариев.

Таким образом, разработка собственной платформы обоснована не отсутствием инструментов автоматизации как таковых, а отсутствием подходящего легковесного решения, объединяющего Architecture-as-Code, inventory, drift detection и reconcile в единую платформенную модель.

\subsection{Постановка задачи разработки платформы}

С прикладной точки зрения задача состоит в том, чтобы дать эксплуатационной команде инструмент регулярной проверки соответствия между архитектурным описанием и фактической средой исполнения. Такой инструмент должен помогать инженеру видеть, какие инфраструктурные workload ожидаются на управляемых хостах, какие из них реально наблюдаются, какие расхождения требуют внимания и какие восстановительные действия могут быть выполнены автоматически в пределах заранее заданных правил.

С учетом проведенного анализа задача выпускной квалификационной работы формулируется следующим образом: разработать и обосновать автоматизированную платформу управления инфраструктурой с поддержкой подхода Architecture-as-Code, которая извлекает desired state из архитектурного описания, получает actual state из источников inventory, обнаруживает drift между ними и инициирует reconcile для поддержанных workload.

Для достижения этой цели необходимо решить следующие инженерные задачи:
\begin{enumerate}
    \item определить модель desired state на основе Structurizr JSON и C4/deployment-представления;
    \item реализовать механизм получения actual state управляемых хостов;
    \item спроектировать сервис обнаружения drift с безопасной обработкой partial-данных;
    \item реализовать Reconciler с операторной моделью и асинхронным выполнением задач;
    \item определить API-контракты между сервисами;
    \item проверить работоспособность платформы тестами и экспериментальной оценкой ключевых сценариев.
\end{enumerate}

Таким образом, первая глава показывает, что рассматриваемая тема находится на пересечении Infrastructure-as-Code, Architecture-as-Code и control-loop автоматизации. Существующие решения дают важные ориентиры, прежде всего Kubernetes-модель controllers/reconcile, однако не закрывают в полной мере задачу легковесного управления host-level workload на основе архитектурного описания. Это обосновывает переход к проектированию собственной платформы, которому посвящена следующая глава.
